% Part: sets-functions-relations
% Chapter: relations
% Section: equivalence-relations

\documentclass[../../../include/open-logic-section]{subfiles}

\begin{document}

\olfileid[pt]{sfr}{rel}{eqv}

\olsection{Relações de Equivalência}

A relação identidade num conjunto é reflexiva, simétrica e
transitiva. Relações~$R$ que têm todas as três destas propriedades são muito
comuns.

\begin{defn}[Relação de equivalência] 
Uma relação $R \subseteq A^2$ que é reflexiva, simétrica e
transitiva chama-se uma \emph{relação de equivalência}. !!^{element}s $x$
e $y$ de~$A$ dizem-se \emph{$R$-equivalentes} se~$Rxy$.
\end{defn}

As relações de equivalência dão origem à noção de \emph{classe de
equivalência}. Uma relação de equivalência ``reparte'' o domínio em
partições distintas. Dentro de cada partição, todos os objetos estão
relacionados uns com os outros; e objetos de partições diferentes não se
relacionam entre si. Por vezes, é útil falar destas
partições \emph{diretamente}. Para esse fim, introduzimos uma
definição:

\begin{defn}\ollabel{def:equivalenceclass}
Seja $R \subseteq A^2$ uma relação de equivalência. Para cada $x \in A$, a
\emph{classe de equivalência} de $x$ em~$A$ é o conjunto $\equivrep{x}{R}
= \Setabs{y \in A}{Rxy}$. O \emph{quociente} de $A$ por~$R$ é
$\equivclass{A}{R} = \Setabs{\equivrep{x}{R}}{x \in A}$, O.s.\@, o conjunto
destas classes de equivalência. 
\end{defn}

O resultado seguinte justifica a definição de classe de equivalência,
provando que as classes de equivalência são de fato as partições de~$A$:

\begin{prop}
Se $R \subseteq A^2$ é uma relação de equivalência, então $Rxy$ se e só se
$\equivrep{x}{R} = \equivrep{y}{R}$.
\end{prop}

\begin{proof}
Para o sentido da esquerda para a direita, suponhamos $Rxy$, e seja $z \in
\equivrep{x}{R}$. Por definição, então, $Rxz$. Como $R$ é uma
relação de equivalência, $Ryz$. (Explicando melhor: como $Rxy$ e~$R$ é
simétrica temos $Ryx$, e como $Rxz$ e~$R$ é transitiva temos~$Ryz$.) Logo $z \in \equivrep{y}{R}$. Generalizando,
$\equivrep{x}{R} \subseteq \equivrep{y}{R}$. Mas de modo perfeitamente análogo,
$\equivrep{y}{R} \subseteq \equivrep{x}{R}$. Logo $\equivrep{x}{R} =
\equivrep{y}{R}$, pela extensionalidade.

Para o sentido da direita para a esquerda, suponhamos $\equivrep{x}{R} =
\equivrep{y}{R}$. Como $R$ é reflexiva, $Ryy$, logo $y \in
\equivrep{y}{R}$. Assim também $y \in \equivrep{x}{R}$ pela suposição
de que $\equivrep{x}{R} = \equivrep{y}{R}$. Logo $Rxy$.
\end{proof}

\begin{ex}
Um bom exemplo de relações de equivalência vem da aritmética modular.
Para quaisquer $a$, $b$, e $n \in \PosInt$, digamos que $a \equiv_n b$ se e só se
dividir $a$ por~$n$ dá o mesmo resto que dividir $b$ por~$n$.
(De modo um pouco mais simbólico: $a \equiv_n b$ se e só se, para algum $k \in
\Int$, $a - b = kn$.) Agora, $\equiv_n$ é uma relação de equivalência, para
qualquer~$n$. E há exatamente $n$ classes de equivalência distintas
geradas por~$\equiv_n$; isto é, $\equivclass{\Nat}{\equiv_n}$ tem
$n$ !!{element}s. São: o conjunto dos números divisíveis por $n$
sem resto, O.s.\@, $\equivrep{0}{\equiv_n}$; o conjunto dos números
divisíveis por $n$ com resto~$1$, O.s.\@, $\equivrep{1}{\equiv_n}$;
\ldots; e o conjunto dos números divisíveis por~$n$ com resto~$n-1$,
O.s.\@,~$\equivrep{n-1}{\equiv_n}$.
\end{ex}

\begin{prob}
Mostre que $\equiv_n$ é uma relação de equivalência, para qualquer $n \in
\PosInt$, e que $\equivclass{\Nat}{\equiv_n}$ tem exatamente $n$ membros.
\end{prob}

\end{document}
